\section{Experimental Setup}
\label{sec:experimental-setup}

We evaluate 15 scanner configurations spanning three categories:
traditional SAST tools, LLM-based scanners, and one hybrid system.
All scanners are run against identical pinned commits of the 26
repositories in the RealVuln corpus. This section describes the scanner
selection, LLM evaluation modes, and reproducibility measures.

\subsection{Scanners Evaluated}
\label{sec:scanners}

\paragraph{Traditional SAST (3 tools).}
We select three widely deployed static analyzers that represent the
dominant analysis strategies in the Python ecosystem:

\begin{itemize}
  \item \textbf{Semgrep} (open-source, rule-based). We invoke
    \texttt{semgrep -{}-config auto} with the default community rule set.
    Semgrep performs intra-procedural pattern matching with limited
    dataflow tracking and is the most commonly adopted open-source SAST
    tool for Python~\cite{semgrep}.
  \item \textbf{Snyk Code} (commercial, pattern + dataflow). We invoke
    \texttt{snyk code test} using the free-tier CLI. Snyk combines
    pattern matching with inter-procedural dataflow analysis and
    machine-learned ranking~\cite{snyk}.
  \item \textbf{SonarQube} (community edition, pattern-based). We use
    the default Python quality profile. SonarQube applies
    intra-procedural taint analysis with an extensive built-in rule
    library~\cite{sonarqube}.
\end{itemize}

\paragraph{LLM-based scanners (11 configurations, 10 distinct models).}
We evaluate ten large language models across eight provider families.
All models receive the same shared prompt (Section~\ref{sec:llm-modes})
and produce structured JSON output. The configurations are:

\begin{itemize}
  \item \textbf{Claude family} (Anthropic): Haiku~4.5
    (single-turn + agentic), Sonnet~4.6 (agentic), Opus~4.6 (agentic).
  \item \textbf{Gemini~3.1 Pro} (Google DeepMind): agentic mode.
  \item \textbf{Grok~3} and \textbf{Grok~4.20 Reasoning} (xAI): both
    agentic mode.
  \item \textbf{Kimi~K2.5} (Moonshot AI): agentic mode.
  \item \textbf{GLM-5} (Zhipu AI): agentic mode.
  \item \textbf{Minimax~M2.7} (MiniMax): agentic mode.
  \item \textbf{Qwen~3.5 397B} (Alibaba Cloud): agentic mode.
\end{itemize}

\paragraph{Hybrid scanner (1).}
\textbf{Kolega v0.0.1} is an LLM-augmented hybrid scanner whose scan
results are published in the repository. Its internals are proprietary;
we include it as a reference point for systems that combine static
analysis with LLM reasoning.

\medskip
\noindent
Table~\ref{tab:scanners} summarizes all evaluated scanners.

\begin{table}[t]
  \caption{Scanner classification. ``Mode'' indicates the primary
    analysis technique. All scanners are run on identical pinned commits.}
  \label{tab:scanners}
  \centering
  \small
  \begin{tabular}{@{}llll@{}}
    \toprule
    \textbf{Scanner} & \textbf{Category} & \textbf{Mode} & \textbf{Open Source} \\
    \midrule
    Semgrep          & Trad.\ SAST & Rule-based        & Yes \\
    Snyk             & Trad.\ SAST & Pattern + dataflow & No (free tier) \\
    SonarQube        & Trad.\ SAST & Pattern-based     & Yes (community) \\
    \midrule
    Claude Haiku 4.5       & LLM & Single-turn + Agentic & No \\
    Claude Sonnet 4.6      & LLM & Agentic               & No \\
    Claude Opus 4.6        & LLM & Agentic               & No \\
    Gemini 3.1 Pro         & LLM & Agentic               & No \\
    Grok 3                 & LLM & Agentic               & No \\
    Grok 4.20 Reasoning    & LLM & Agentic               & No \\
    Kimi K2.5              & LLM & Agentic               & No \\
    GLM-5                  & LLM & Agentic               & No \\
    Minimax M2.7           & LLM & Agentic               & No \\
    Qwen 3.5 397B          & LLM & Agentic               & No \\
    \midrule
    Kolega v0.0.1    & Hybrid      & Proprietary       & No \\
    \bottomrule
  \end{tabular}
\end{table}

\subsection{LLM Evaluation Modes}
\label{sec:llm-modes}

We evaluate LLM-based scanners in two modes that represent different
points on the cost--capability spectrum.

\paragraph{Single-turn (pilot).}
The entire repository context is concatenated into a single API call.
The model receives the source files and the shared prompt, and must
return all findings in one response. This is the cheapest mode---it
requires exactly one model invocation per repository---but it is
constrained by the model's context window and receives no feedback
on its intermediate reasoning. We use this mode as a lower-bound
baseline; only Claude Haiku~4.5 is evaluated in single-turn mode (in
addition to its agentic evaluation).

\paragraph{Agentic.}
The model operates within a tool-use loop, with access to file reading,
\texttt{grep}-style search, and shell execution tools. It can
iteratively explore the codebase, form hypotheses about data flow, and
refine its findings across multiple turns. This mode is more expensive
(typically 10--50$\times$ more tokens than single-turn) but more
closely mirrors how a human auditor would approach the task. All 10
LLM models are evaluated in agentic mode.

\paragraph{Shared prompt.}
All LLM configurations receive the same prompt template, identified by
its content hash (\texttt{sha256:828b00245b42}). The prompt instructs
the model to act as a security auditor, specifies the expected JSON
output schema (CWE identifier, file path, line number, severity), and
provides minimal guidance on vulnerability categories. Per-model prompt
tuning is intentionally avoided: while a shared prompt may
advantage or disadvantage particular models, it ensures that differences
in performance reflect model capability rather than prompt engineering
effort. The full prompt is published in the benchmark repository.

\paragraph{Output schema.}
Each LLM must return a JSON array of findings, where each finding
contains: a CWE identifier (\texttt{cwe}), the affected file path
(\texttt{file}), a line number (\texttt{line}), and a severity label
(\texttt{severity}). Findings that do not conform to this schema are
discarded during parsing.

\paragraph{Cost tracking.}
We record input and output token counts for every LLM invocation and
compute per-scan costs using provider pricing as of March~2026. This
enables cost-efficiency analysis (e.g., F3 score per dollar) in
Section~\ref{sec:results}, which is critical for practitioners
evaluating deployment trade-offs.

\subsection{Reproducibility}
\label{sec:reproducibility}

All artefacts needed to reproduce or extend the evaluation are committed
to the public repository:

\begin{itemize}
  \item \textbf{Scan results.} Raw scanner output for every
    (repository, scanner) pair is stored in
    \texttt{scan-results/\{repo\}/\{scanner\}/results.json}.
  \item \textbf{Version locking.} The file
    \texttt{benchmark-manifest.json} records the content hash of each
    ground-truth file, the pinned commit SHA for every target
    repository, and the prompt version hash, ensuring that scores can be
    tied to a specific snapshot of the benchmark.
  \item \textbf{Re-scoring.} Any user can re-derive all metrics by
    running:
    \begin{verbatim}
  python score.py --repo <repo> --all-scanners
    \end{verbatim}
  \item \textbf{Dashboard regeneration.} The interactive HTML dashboard
    is regenerated via:
    \begin{verbatim}
  python dashboard.py --scanner-group all
    \end{verbatim}
\end{itemize}

\noindent
Together, these measures ensure that the evaluation is fully
transparent: readers can verify every number in this paper, and future
researchers can add new scanners or repositories without modifying the
scoring pipeline.
